\chapter{Future Work}

The current system has successfully achieved its main goals, including cross-platform integration, AI-powered evaluation, semantic search, and enhanced agent search capabilities (including multi-region support, gradual filter relaxation, fallback explanations, and evaluation chat assistants). However, there is still room for improvement in future enhancements.

\section{Data Pipeline and Feature Enrichment}

\subsection{Geospatial and Proximity Features}
Our current process uses the Singapore Land Authority's \textbf{OneMap API} for address geocoding and reverse geocoding to generate district codes. However, our rating model is primarily based on building characteristics (living area, number of bedrooms, use, age of the building, and neighborhood). In future versions, we plan to expand OneMap integration to calculate accurate walking routes to important facilities such as MRT stations, top primary schools, and the Central Business District (CBD), and integrate this as an additional model feature.

The expected performance improvement can be quantified by the error analysis in Section~\ref{sec:error_analysis}. The main reason for the residual error of the apartment sales model (current MAPE 13.35\%) is the \textit{micro-location premium in the area}. This premium can be accurately reflected by the distance to the MRT station and the distance to the school. A study on the Singapore hedonistic model showed that the MAPE of the apartment segment could be reduced by 3--5 percentage points after considering the proximity factor \cite{park2015using}. A conservative estimate of 3\,pp would reduce the MAPE of apartment sales from 13.35\% to about \textbf{10\%}, which would be close to the accuracy of the current Housing and Development Board (HDB) sales model. For the landed sales segment, the MAPE is expected to be reduced from 26.28\% to 18--20\% due to the distance to amenities at the plot level. This implementation requires one-time batch geocoding of 53,497 merged records using the OneMap Routing API (free tier, 250 requests/min), followed by feature engineering and model retraining---an estimated two-week effort.

\subsection{Expanded Team Coverage}
While the PropertyGuru scraper developed by the author is fully operational and covers the majority of the market, the team plans to expand the \texttt{condo\_basic} master records (managed by other team members) to include historical transaction data explicitly linked to current listings.

\section{Advanced AI Modeling}

\subsection{Domain-Adapted Language Model for Property Search}
The current semantic search pipeline uses a universal language model (Claude) with system requests that specify the context of the Singaporean real estate market. A promising approach for the future would be to fine-tune a smaller, open-source language model (e.g., Llama 3 or Mistral) based on a curated dataset of search queries related to the Singaporean real estate market and the corresponding structured filters. This finely tuned model can improve accuracy while reducing latency and API costs for domain-specific representations such as HDB flat types and postcodes in Singapore.

\subsection{Computer Vision for Condition Assessment}
Real estate information often contains a wealth of visual information about the quality of interior finishes and natural lighting. However, this information cannot currently be quantified using tabular models. Therefore, we propose using a pre-trained Vision Transformer (ViT) or a Convolutional Neural Network (CNN) to extract aesthetic elements from real estate images. This allows the valuation model to differentiate between "unrenovated" and "designer-renovated" properties of similar size and age.

\subsection{Spatio-Temporal Market Forecasting}
The current gradient-boosting model predict valuations at a specific point in time based on current real estate information. By collecting long-term datasets of historical transactions and expired real estate information in the coming months, we aim to develop a time-series forecasting model (e.g., using recurrent neural networks or spatiotemporal graph neural networks). This model will enable the system to predict future asset appreciation and rental yield trends.

\subsection{Automated Listing Freshness and Alert System}
The current system only scrapes listings when it needs to. In the future, regularly updated crawlers will detect changes and notify registered users when the prices of offers matching their saved search criteria change, or when offers are removed or added. This changes the platform from a tool for passive searching to a service for actively monitoring the market.

\section{Evaluation and Validation Improvements}

\subsection[Larger-Scale External Benchmark for NL Search]{Larger-Scale External Benchmark for Natural Language Search}
\label{sec:future_blind_eval}

Three annotators conducted independent blind evaluations of 18 queries (Section~\ref{sec:blind_eval}). 
The evaluation found no outright errors, and the adjusted F1 score was approximately 91\%. 
The next step in the validation process is to conduct a more thorough external benchmark. 
This formal evaluation uses at least $N \geq 100$ queries. These are created by annotators 
completely independent of the project team. It includes a wider range of attribute types, 
edge constraints, and idiomatic expressions of Singaporean English (e.g., ``5I flat'', 
``DBSS'', ``mickey mouse unit''). External benchmarks improve the statistical power of 
error pattern recognition. They also lead to more robust F1 confidence intervals. 
Ultimately, these reviews make the system a more reliable search interface.

\subsection{HDB Lease Risk Badge}
\label{sec:future_lease_badge}

The feature ablation study (Section~\ref{sec:policy_signal}) showed that
\texttt{property\_age} carries the second-largest SHAP contribution in the HDB Sale
model, and its removal causes a 130.1\,pp MAPE degradation---a sensitivity consistent
with Singapore's CPF lease-eligibility policy creating a demand-side price
discontinuity for ageing HDB stock. This behaviour has been operationalised into a
user-facing \textit{lease risk badge} displayed on every HDB listing detail page (see
Section~\ref{sec:lease_badge_impl} of the Implementation chapter for full details).
The badge classifies each HDB listing into one of three CPF eligibility tiers based on
remaining lease duration, making the model's age-sensitivity explicitly actionable for buyers.
Future work could extend this to a proactive alert system that flags listings approaching
a tier boundary within the next three years.

\section{Recommendation Engine Extensions}

\subsection{Collaborative Filtering at Scale}

The current recommendation engine is a hybrid content-based and valuation-grounded ranker,
chosen because the platform is at the cold-start stage: fewer than ten registered users with
sparse interaction histories make collaborative filtering (CF) unreliable. As the user base
grows beyond a few hundred active users with at least 5--10 saves each, matrix factorisation
approaches---such as the implicit-feedback ALS model of Hu, Koren, and
Volinsky~\cite{hu2008collaborative} or Neural Collaborative Filtering~\cite{he2017neural}---become
viable as a second-stage re-ranker. The latent user and item embeddings learnt by CF would
capture cross-user preference patterns that are invisible to the current single-user
profile approach: for example, users who save the same Condo D9 listings are likely to
share a broader preference cluster that the content-based component cannot infer from
attributes alone.

A practical migration path is a \textit{cascade hybrid}~\cite{burke2002hybrid}: the existing
content-based + bargain-score engine generates a top-200 candidate set (as it does today),
and a CF model re-ranks the top-200 using learned user--item affinities. This preserves
the cold-start resilience of the content-based layer while layering CF's generalisation
capability on top once sufficient interaction data exists.

\subsection{Implicit Feedback Signals}

The current system records only \textit{explicit} positive feedback (saves). Implicit
signals available in the frontend---listing click-throughs, map viewport dwells, detail
page scroll depth, and chat assistant queries on a specific listing---constitute a much
richer and more frequent interaction stream. Incorporating these as pseudo-positive
feedback (with lower confidence weight than saves) would enable the preference profile
to update in near real-time during a browsing session, and would allow the engine to
bootstrap recommendations for users who have not yet saved any listings.

\subsection{Contextual and Online Learning}

Both the content weights ($w_{\text{type}}, w_{\text{district}}, \ldots$) and the
bargain-score threshold are currently fixed constants tuned by offline evaluation.
A contextual bandit formulation~\cite{li2010contextual} would allow the system to
adaptively tune these weights per user based on observed click-through rates on
recommended listings, converging towards a personalised weight vector for each user
over time. This is a natural extension once A/B testing infrastructure is in place.
